\section{Discussion}
\label{sec:discussion}

\subsection{Trade-off analysis}
The experiments let us reason about \emph{when} a technique earns its cost.

\paragraph{Accuracy vs.\ latency.}
Re-ranking buys a real ordering improvement (MRR $0.866\!\to\!0.876$, Context
Precision $0.655\!\to\!0.680$) for only ${\sim}0.4$\,s and is clearly worth it.
Query enhancement is the opposite: it costs ${\sim}1.8$\,s \emph{and} lowers
every quality metric, so it is on the wrong side of every trade-off and is
dropped. The recommended pipeline (stage 05) keeps end-to-end latency around
$3.6$\,s, within the project's $<5$\,s target for text queries.

\paragraph{Complexity vs.\ gain.}
Intent routing adds an LLM call and code complexity but does not change retrieval
quality. Its justification is not accuracy but \emph{safety and correctness}: it
delivers a 100\% redirect rate on pricing and out-of-scope queries, structurally
preventing the chatbot from inventing prices --- a hard business requirement.
This is complexity that pays for itself in guarantees rather than metrics.

\paragraph{Model size vs.\ quality.}
The embedding study shows the largest model is not required: the compact,
domain-matched \texttt{bkai} encoder ties the larger \texttt{e5-base}, while the
even smaller \texttt{MiniLM} is the one that genuinely hurts retrieval. The
minimum viable embedder for this corpus is a Vietnamese-aware model of
${\sim}135$M parameters.

\paragraph{The headline lesson.}
\Cref{fig:cumulative} (RAGAS view) and \Cref{fig:ladder} (IR view) together make
the project's central point: ``more advanced'' is not monotonically better. The
best system is obtained by adding hybrid retrieval and re-ranking, adding intent
routing for safety, and \emph{deliberately omitting} query enhancement. Knowing
which optimization to leave out is itself a result.

\begin{figure}[H]
\centering
\includegraphics[width=0.8\textwidth]{cumulative_gains.png}
\caption{RAGAS metric view of the cumulative ladder, complementing the IR view
in \Cref{fig:ladder}.}
\label{fig:cumulative}
\end{figure}

\subsection{Threats to validity}
\begin{itemize}[leftmargin=1.4em,itemsep=2pt]
  \item \textbf{Test-set size.} 80 questions is larger than our initial draft but
        still modest; the LLM-judged RAGAS faithfulness mean in particular has
        run-to-run variance of ${\pm}0.05$--$0.08$. We mitigate this by leading
        with deterministic IR metrics, which are stable across runs, and by
        treating sub-$0.05$ faithfulness differences as noise.
  \item \textbf{Single judge.} RAGAS uses a single Gemini Flash-Lite judge; a
        stronger or ensemble judge would reduce noise but increase cost. The
        relative rankings (especially of retrieval techniques) are corroborated
        by the judge-independent IR metrics.
  \item \textbf{Reference-based metrics.} Faithfulness and context recall depend
        on the reference answers; we rebuilt the test set directly from the
        source data to remove the fabricated-reference problem present in an
        earlier draft (\Cref{sec:testset}).
  \item \textbf{Not evaluated.} The Qwen/Ollama LLM comparison and the CLIP
        image-matching accuracy were not run (no local Ollama server; no held-out
        test images). The image path is implemented but its accuracy is future
        work.
\end{itemize}

\section{Conclusion}
\label{sec:conclusion}

We built an evaluation-driven Vietnamese RAG chatbot for a printing-and-packaging
catalog and used it to study how each RAG component contributes to quality. By
adding one technique at a time and measuring with both RAGAS and classical IR
metrics on an 80-question test set, we established a clear, defensible result:

\begin{enumerate}[leftmargin=1.5em,itemsep=2pt]
  \item \textbf{Hybrid retrieval + cross-encoder re-ranking} is the quality win
        (Recall@5 $0.884\!\to\!0.899$, MRR $0.821\!\to\!0.876$, Context Precision
        $\to0.680$) for negligible added latency.
  \item \textbf{Intent routing} adds a 100\% pricing/out-of-scope redirect rate
        --- a business-critical guarantee --- at no retrieval cost.
  \item \textbf{Query enhancement (HyDE/rewrite)} is rejected: on a small,
        on-vocabulary catalog it causes vocabulary drift that collapses recall
        and adds latency.
  \item \textbf{Vietnamese-aware embeddings} (\texttt{bkai} or \texttt{e5-base})
        clearly beat a generic small multilingual model; chunking strategy is
        immaterial at this scale.
\end{enumerate}

The recommended production pipeline is hybrid retrieval + re-ranking + intent
routing with the \texttt{bkai} embedder and document-level chunking. Beyond the
specific configuration, the project demonstrates the methodological point that
the most elaborate pipeline is not the best one --- and that disciplined,
component-wise evaluation is what reveals which optimizations to keep and which
to discard. Future work includes the open-source LLM comparison, a quantitative
evaluation of the CLIP image-matching path, and a hybrid-$\alpha$ sweep.
